\section{Cross Correlation (Series Dot Product)}

The cross correlation is a measure of the {\em similarity  
in time}
of two time series signals.  It is often referred to as
the {\em dot product} of two time series because 
the cross correlation operation is a dot production operation.

\subsection{Vector Dot Product}
Consider two vectors $\va$ and $\vb$ in three-dimensional
space, measured from the same origin and with the same 
coordinate system, with components,
\begin{align}
\va & = \left[ \; a_1, a_2, a_3 \; \right], \\
\vb & = \left[ \; b_1, b_2, b_3 \; \right]. \\
\end{align}
The dot product of the two vectors is simply,
\begin{align}
 \va \cdot \vb = a_i \; b_i 
 & = a_1 \; b_1 + a_2 \; b_2 + a_3 \; b_3 \\
 \mbox{or generally} 
 & = 
 \sum_{i=1}^n a_i \; b_i.
\end{align}

Next, consider $\va(t)$ and $\vb(t)$ be two time-varying
signals, occurring over time interval $[\;t_0, t_f\;]$. 
Let the amplitudes of $\va(t)$ 
and $\vb(t)$ be measured from zero to $t_f$ at $n$ equally
distributed points such that the signals are sampled 
at $n$ uniform discrete times in the interval.
The two vectors are written explicitly as
\begin{align}
\va & = 
\left[ 
\; a_0, a_1, a_2, \ldots, 
 a_{n-3},
 a_{n-2},
 a_{n-1}
 \;
\right],  \\
\vb & = 
\left[ 
\; b_0, b_1, b_2, \ldots, 
 b_{n-3},
 b_{n-2},
 b_{n-1}
 \;
\right].
\end{align}

\begin{Hremark}
To account for the {\em initial condition} of the time series, 
we decrement the counter, 
\be
1 \ldots n,
\ee
to become 
\be
0 \ldots n-1.
\ee
In both of these cases, the number of sample 
points is identical, but the indexing is different by one (1). 
\end{Hremark}

\begin{Hremark}
Notice that the two signals 
$\va(t)$ and $\vb(t)$ are {\em synchronized} in time
such that 
\be
\vt= \left[ 
\; t_0, t_1, t_2, \ldots,
 t_{n-3},
 t_{n-2},
 t_{n-1}
 \;
\right].
\ee
Notice that $\Delta t = t_1 - t_0 = t_2 - t_1 = \cdots
t_{n-1} - t_{n-2}$ and that the signals
are sampled at a frequency 
of $\frac{1}{\Delta t}$~Hz.
\end{Hremark}

\subsection{Series Dot Product}
The dot product of these two signals is
\begin{align}
 \va \cdot \vb = a_i \; b_i 
 & = a_0 \; b_0 + a_1 \; b_1 + a_2 \; b_2 + \cdots + 
 a_{n-1} \; b_{n-1} \\
 \mbox{or generally} 
 & = 
 \sum_{i=0}^{n-1} a_i \; b_i.
\end{align}
For comparison of two time series, the dot product will be referred 
to as the {\bf cross correlation score}.

\begin{Hexample}
	{
    This example is a reproduction of the cross correlation example
	from https://anomaly.io/understand-auto-cross-correlation-normalized-shift/.
	Let vectors $\va$, $\vb$, and $\vc$ be defined as follows,
	\begin{align}
		\va &= [\; 1, 2, -2, 4, 2, 3, 1, 0 \;] \\
		\vb &= [\; 2, 3, -2, 3, 2, 4, 1, -1 \; ] \\
		\vc &= [\; -2, 0, 4, 0, 1, 1, 0, -2 \;]
	\end{align}
	Then the cross-correlation score between $\va$ and $\vb$ is found to be 
	high, 
	\be
	\va \cdot \vb = 41; 
	\ee
	whereas, the cross-correlation score between $\va$ and $\vc$ is found to 
	be low, 
	\be
	\va \cdot \vc = -5.
	\ee
	}
\end{Hexample}

\subsection{Limitation of Cross Correlation}
A significant limitation of the cross correlation method is its 
dependence on {\em amplitude}:
\begin{itemize}
\item Large postive magnitude increase the cross correlation.
\item Large negative magnitude decrease the cross correlation.
\item A cross correlation score shouldn't get better (positive) or
worse (negative) if the amplitude of the signal is increased. 
\item Shape similarity is obfuscasted by amplitude dependence.  For 
example, the three cross correlation scores generated from comparison 
of three signals, $A \sin(t)$, $B \sin(t)$, and $C \sin(t)$, 
all of the same shape, $\sin(t)$, but with unique
amplitudes $A, B, C$, will report that two out of the three signals are 
better correlated by virtue of their closer amplitudes.  In a sense, 
this result is satisfactory, since amplitude distinguishes one signal
from another.  However, this result is incomplete, since the amplitude 
dependence obfuscastes the fact that all three signals have the same 
basic shape.
\item Shape similarity obfuscation is easily seen with shape 
self-similarity, that is, using the same signal for both vectors in the
cross correlation.  In this case 
\be
  \va \cdot \va > \va \cdot \frac{\va}{2}, 
\ee
which is a sensible outcome, but not helpful to distinguish shape
similarity.
\end{itemize}

To remedy the shortcoming of cross correlation caused by amplitude
dependence, {\bf vector normalization} is employed, as explained 
in a following section.  Prior to this, however, we make a brief 
introduction to the p-norm.

\subsection{p-norm}
For any any number $p \in \real$ and $p \geq 1$, the {\bf p-norm} 
is defined as
\be
 \parallel \va \parallel_p
 \; \defe \left( 
 \sum_{i=0}^{n-1} \mid a_i \mid^p 	 
 \right)^{\frac{1}{p}}.
\ee

\begin{Hexample}
	{
		The L2-norm, also called the {\bf Euclidean norm}, is defined as 
		\be
		 \parallel \va \parallel_2
		 \; \defe \left( 
		 \sum_{i=0}^{n-1} \mid a_i \mid^2 	 
		 \right)^{\frac{1}{2}},
		\ee
		and has the physical interpretation as the length of the vector
		on which the L2-norm operates.
	
		\medskip

		The L1-norm, also called the {\bf Manhattan taxicab norm}, is 
		defined as 
		\be
		 \parallel \va \parallel_1
		 \; \defe \left( 
		 \sum_{i=0}^{n-1} \mid a_i \mid 	 
		 \right),
		\ee
		and has the physical interpretation as the cumulative length 
		(mileage) accumulated for each segement $a_i$.

		\medskip

		The {\bf absolute value} is the L1-norm on a one-dimensional
		vector space, 
		\be
		 \parallel \va \parallel_1
		 \; \defe \left( 
		 \sum_{i=0}^{0} \mid a_i \mid 	 
		 \right) 
		 = 
		 \; \mid a_0 \mid,
		\ee
	}
\end{Hexample}

\subsection{Normalized Cross Correlation}
When vector normalization is used with cross correlation, the resulting
method is referred to as {\em normalized cross correlation}.  
Normalized cross correlation is the same as cross correlation with 
one exception, the vectors used in the dot product are normalized by 
their L2-norm.  

The normalized cross correlation for the three-dimensional vectors
$\va$ and $\vb$ is defined as 
\begin{align}
\frac{\va}{\parallel \va \parallel}_2
\cdot
\frac{\vb}{\parallel \vb \parallel}_2
& \defe 
\frac{a_i}{\sqrt{\sum_{i=1}^n a_i \; a_i}} 
\cdot
\frac{b_i}{\sqrt{\sum_{i=1}^n b_i \; b_i}}  \\
& = 
\frac{a_1 \; b_1 + a_2 \; b_2 + a_3 \; b_3}{
    \left(\sqrt{a_1^2 + a_2^2 + a_3^2}\right) \; 
    \left(\sqrt{b_1^2 + b_2^2 + b_3^2}\right)}. 
\end{align}
The foregoing result generalizes to time series $\va(t)$ and $\vb(t)$ 
as 
\begin{align}
\frac{\va}{\parallel \va \parallel} 
\cdot
\frac{\vb}{\parallel \vb \parallel} 
& \defe 
  \frac{\dst\sum_{i=0}^{n-1} a_i \; b_i}
  {\left(\sqrt{\dst\sum_{i=0}^{n-1} a_i^2}\right) \; 
   \left(\sqrt{\dst\sum_{i=0}^{n-2} b_i^2}\right)} 
\end{align}

Consider again the vectors $\va$, $\vb$, and $\vc$ from 
Example~1.

\begin{Hexample}
	{
    Consider again the vectors $\va$, $\vb$, and $\vc$ from 
    Example~1.
	The {\em normalized} cross-correlation score between $\va$ and $\vb$ is found to be 
	high (relative to a maximum of 1.0), 
	\be
    \frac{\va}{\parallel \va \parallel} \cdot 
    \frac{\vb}{\parallel \vb \parallel} = \frac{41}{\sqrt{39} \; \sqrt{48}} 
    = 0.947;
	\ee
    whereas, the {\em normalized} cross-correlation score between $\va$ and $\vc$ 
    is found to 
	be low (relative to no correlation for 0.0), 
	\be
    \frac{\va}{\parallel \va \parallel}
    \cdot 
    \frac{\vc}{\parallel \vc \parallel} = \frac{-5}{\sqrt{39} \; \sqrt{26}} = -0.157.
	\ee
	}
\end{Hexample}

The {\bf benefits} of {\em normalized} cross-correlation are two-fold:  
Amplitude independence and unit normal scoring.  Amplitude independence 
is seen as 
\be
\va \cdot \va = \va \cdot \frac{\va}{2}.
\ee
The normalized cross-correlation score now ranges betweeen -1 and +1.
\begin{itemize}
 \item +1 is perfect correlation,
 \item 0 is no correlation, and
 \item -1 is anti-correlation (correlated, but in opposite phase).
\end{itemize}

\subsection{Discrete Case in Matrix Form}
In the discrete sense, 
\be
Y_i = W_{ij} X_j,
\ee
which takes the explicit form,
\be
\begin{tabular}{c}
slide- \\
forward \\
operation \\ 
$\vdots$ \\
len = $n$\\
\hline
slide- \\ 
backward \\
operation \\
$\vdots$ \\
len = $n-1$
\end{tabular}
\hspace{0.5cm}
%
\left\{
\begin{tabular}{c}
$y_0$ \\
$y_1$ \\
$y_2$ \\ 
$\vdots$ \\
$y_{n-1}$ \\
\hline
$y_{-1}$ \\
$y_{-2}$ \\
$y_{-3}$ \\
$\vdots$ \\
$y_{-(n-1)}$ 
\end{tabular}
\right\}_{(2n-1 \times 1)}
= \nonumber
\ee 
\be
\left[
\begin{tabular}{ccccccc}
$w_0$ & 
$w_1$ & 
$w_2$ & 
$\cdots$ & 
$w_{n-3}$ & 
$w_{n-2}$ & 
$w_{n-1}$  
 \\
0 & 
$w_0$ & 
$w_1$ & 
$w_2$ & 
$\cdots$ & 
$w_{n-3}$ & 
$w_{n-2}$ 
 \\
0 & 
0 & 
$w_0$ & 
$w_1$ & 
$w_2$ & 
$\cdots$ & 
$w_{n-3}$ 
 \\
 & 
 & 
 & 
$\ddots$ & 
 & 
 & 
$\vdots$ 
\\
 & 
 & 
 & 
 & 
$\ddots$ & 
 & 
$\vdots$ 
 \\
 & 
 & 
 & 
 & 
 & 
$\ddots$ & 
$\vdots$ 
 \\
0 & 
0 & 
0 & 
0  & 
0 & 
$\cdots$ & 
$w_0$ 
 \\ \hline
$w_1$ & 
$w_2$ & 
$\cdots$ & 
$w_{n-3}$ & 
$w_{n-2}$ & 
$w_{n-1}$ & 
0 
\\
$w_2$ & 
$\cdots$ & 
$w_{n-3}$ & 
$w_{n-2}$ & 
$w_{n-1}$ & 
0 & 
0
\\
$\vdots$ & 
$w_{n-3}$ & 
$w_{n-2}$ & 
$w_{n-1}$ & 
0 & 
0 & 
0 
\\
$\vdots$ & 
 & 
\reflectbox{$\ddots$} & 
 & 
 & 
 & 
\\
$\vdots$ & 
\reflectbox{$\ddots$} & 
 & 
 & 
 & 
 & 
 
\\
$w_{n-1}$ & 
0 & 
0 & 
0 & 
0 & 
0 & 
0 
\end{tabular}
\right]_{(2n-1 \times n)}
\left\{
\begin{tabular}{c}
$x_0$ \\ 
$x_1$ \\ 
$x_2$ \\ 
$\vdots$ \\ 
$x_{n-3}$ \\ 
$x_{n-2}$ \\ 
$x_{n-1}$ \\ 
\end{tabular}
\right\}_{(n \times 1)}
\ee

\subsection{Continuous Case Generalized from the Discrete Case}
Finally, let the cross correlation function $y(t)$ 
be the sliding inner product of the signal $x(t)$ with
the weighting function $w(t)$ such that 
\be
y(t) = \int_0^{t_f} x(t) \; w(t - \tau) \; d\tau,
\ee
which is the filter {\bf slide-forward option}, and 
\be
y(t) = \int_{-t_f}^{0} x(t) \; w(t - \tau) \; d\tau 
\ee
is the {\bf slide-backward operation}. 
The two forgoing operations
can be combined to generally be the {\bf 
sliding dot product operation} as 
\be
y(t) = \int_{-t_f}^{t_f} x(t) \; w(t - \tau) \; d\tau.
\ee

